% Polished version of abs-xiao-2.tex (proofread by Claude). The original file is unchanged.
\begin{abstract}
Geometric spatial reasoning underlies applications ranging from robotics to automated design, yet whether language models can turn a verbal description into a physically connected 3D assembly remains largely unexplored. LEGO\textsuperscript{\tiny\textregistered} brick assembly offers a demanding constructive test of this ability: a short, underspecified prompt leaves dimensions, proportions, and part layout implicit, while a valid assembly must both capture the prompt's semantic intent and satisfy the physical constraints imposed by interlocking studs. We formalize this setting as the \texttt{Text-to-Brick} task and present BrickForma, a framework that crystallizes the geometric knowledge of large language models (LLMs) into a small language model for this task without requiring human-annotated or human-verified 3D shape assets. In BrickForma, executable programs written by LLM agents procedurally generate 80,000 voxel models across 21 categories on a $48\times48\times48$ grid. Each model is paired with five captions of increasing detail, yielding approximately 400,000 geometry-grounded training examples. Connectivity analysis and refinement convert the voxel models into geometrically checked brick targets, of which 79,301 pass all checks, and we fine-tune Qwen3-1.7B to predict brick dimensions and positions from text. Prior systems depend on curated assets: BrickGPT builds on human-annotated ShapeNetCore meshes, and BrickNet relies on human-designed LDraw assemblies. BrickForma instead generates both the voxel geometries and their brick targets programmatically. On 525 shared prompts spanning five description levels, BrickForma's raw outputs achieve the highest VQAScore, PE, and SigLIP2 scores of the three systems at every level. Their invalid-placement ratio is 0.79\%, close to BrickGPT's 0.00\% and far below BrickNet's 18.42\%.
\end{abstract}
